% !TEX root = ../dabirestan.tex
%% ============================================================
%%  فصل ۱ — احتمال چیست؟ آزمایش، فضای نمونه و پیشامد
%% ============================================================
\chapter{احتمال چیست؟}
\label{hs:ch1}

\chapterintro{\begin{itemize}
  \item احتمال یعنی چه و چرا آن را با عددی بین ۰ و ۱ نشان می‌دهیم؛
  \item آزمایش تصادفی، فضای نمونه و پیشامد چیست؛
  \item چطور پیشامدها را با «یا»، «و» و «نه» ترکیب کنیم؛
  \item احتمال تجربی را چطور با تکرار آزمایش به دست آوریم.
\end{itemize}}

%% ------------------------------------------------------------
\section{شانس را با عدد بیان کنیم}
%% ------------------------------------------------------------

در زندگی روزمره مدام دربارهٔ «شانس» حرف می‌زنیم: «احتمالاً امروز باران نمی‌آید»، «شانس قبولی‌ام زیاد است»، «برنده شدن در قرعه‌کشی تقریباً غیرممکن است». در ریاضی این شانس‌ها را با یک \emph{عدد بین ۰ و ۱} بیان می‌کنیم و اسمش را \term{احتمال}{probability} می‌گذاریم:

\begin{figblock}
\begin{tikzpicture}[x=13cm]
  \shade[left color=warncolor!35, right color=excolor!40] (0,-0.18) rectangle (1,0.18);
  \draw[primary, thick] (0,-0.18) rectangle (1,0.18);
  \foreach \x/\l/\t in {0/0/{غیرممکن}, 0.25/{\tfrac14}/{بعید}, 0.5/{\tfrac12}/{پنجاه‌پنجاه}, 0.75/{\tfrac34}/{محتمل}, 1/1/{حتمی}} {
    \draw[primary, thick] (\x,-0.28) -- (\x,0.28);
    \node[font=\small] at (\x,-0.62) {\lr{$\l$}};
    \node[font=\small\semifont, text=primary] at (\x,0.65) {\rl{\t}};
  }
  \node[font=\scriptsize, text=graytext, align=center] at (0,1.25) {\rl{خورشید از غرب طلوع کند}};
  \node[font=\scriptsize, text=graytext, align=center] at (0.5,1.25) {\rl{سکه شیر بیاید}};
  \node[font=\scriptsize, text=graytext, align=center] at (1,1.25) {\rl{فردا صبح شود}};
\end{tikzpicture}

\figcaption[hs:fig:scale]{مقیاس احتمال: هر چه احتمال به ۱ نزدیک‌تر باشد، رخ دادن آن پیشامد محتمل‌تر است.}
\end{figblock}

\begin{itemize}
  \item احتمال \textbf{۰} یعنی آن اتفاق \emph{هرگز} نمی‌افتد.
  \item احتمال \textbf{۱} یعنی آن اتفاق \emph{حتماً} می‌افتد.
  \item احتمال $\tfrac12$ یعنی شانس افتادن و نیفتادنش برابر است، مثل شیر آمدن یک سکهٔ سالم.
\end{itemize}

احتمال را گاهی به صورت درصد هم می‌گوییم: احتمال $0.3$ همان ۳۰٪ است. پیش‌بینی هواشناسی «۷۰٪ احتمال بارش» یعنی احتمال باران $0.7$ است.

%% ------------------------------------------------------------
\section{آزمایش تصادفی و فضای نمونه}
%% ------------------------------------------------------------

\begin{definition}[label=hs:def:experiment]{آزمایش تصادفی، برآمد، فضای نمونه}
\begin{itemize}
  \item \term{آزمایش تصادفی}{random experiment} کاری است که نتیجه‌اش را از قبل نمی‌دانیم؛ مثل پرتاب سکه یا تاس.
  \item هر نتیجهٔ ممکن آزمایش را یک \term{برآمد}{outcome} می‌نامیم.
  \item مجموعهٔ \emph{همهٔ} برآمدهای ممکن را \term{فضای نمونه}{sample space} می‌نامیم و با $S$ نشان می‌دهیم. تعداد اعضای $S$ را $n(S)$ می‌نویسیم.
\end{itemize}
\end{definition}

\begin{example}[label=hs:ex:spaces]{چند فضای نمونه}
\begin{enumerate}
  \item پرتاب یک سکه: دو برآمد «شیر» و «خط» داریم. برای کوتاه‌نویسی شیر را $H$ و خط را $T$ می‌نامیم: \ $S=\{H,T\}$ و $n(S)=2$.
  \item پرتاب یک تاس: \ $S=\{1,2,3,4,5,6\}$ و $n(S)=6$.
  \item پرتاب دو سکه (یکی پس از دیگری): \ $S=\{HH,\ HT,\ TH,\ TT\}$ و $n(S)=4$.
  \item انتخاب یک روز هفته به تصادف: \ $n(S)=7$.
\end{enumerate}
\end{example}

برای آزمایش‌هایی که چند مرحله دارند، \emph{نمودار درختی} کمک می‌کند هیچ برآمدی را از قلم نیندازیم:

\begin{figblock}
\begin{tikzpicture}[
  grow=left, level distance=2.6cm,
  level 1/.style={sibling distance=2.2cm},
  level 2/.style={sibling distance=1cm},
  every node/.style={font=\small},
  st/.style={draw=primary!60, fill=primarylight, rounded corners=2pt, inner sep=3pt},
  edge from parent/.style={draw=primary!55, thick}]
  \node[circle, fill=primary, inner sep=2pt] {}
    child { node[st] {\lr{$H$}}
      child { node[st] (hh) {\lr{$H$}} }
      child { node[st] (ht) {\lr{$T$}} } }
    child { node[st] {\lr{$T$}}
      child { node[st] (th) {\lr{$H$}} }
      child { node[st] (tt) {\lr{$T$}} } };
  \foreach \n/\l in {hh/HH, ht/HT, th/TH, tt/TT} \node[anchor=east, font=\small, text=accent!80!black] at ([xshift=-3mm]\n.west) {\lr{$\l$}};
  \node[font=\scriptsize, text=graytext] at (-2.6,2.0) {\rl{سکهٔ اول}};
  \node[font=\scriptsize, text=graytext] at (-5.2,2.0) {\rl{سکهٔ دوم}};
\end{tikzpicture}

\figcaption[hs:fig:tree2]{نمودار درختی پرتاب دو سکه: هر مسیر از ریشه تا انتها یک برآمد است.}
\end{figblock}

\begin{warning}
در پرتاب دو سکه، برآمدهای $HT$ (اولی شیر، دومی خط) و $TH$ (اولی خط، دومی شیر) \emph{دو برآمد متفاوت}‌اند. اگر فکر کنید فقط سه حالت «دو شیر، دو خط، یکی از هر کدام» وجود دارد، در محاسبهٔ احتمال اشتباه می‌کنید (در فصل ۳ دلیلش را می‌بینیم).
\end{warning}

%% ------------------------------------------------------------
\section{پیشامد}
%% ------------------------------------------------------------

\begin{definition}[label=hs:def:event]{پیشامد}
هر زیرمجموعه از فضای نمونه یک \term{پیشامد}{event} نام دارد. اگر برآمد آزمایش عضوی از پیشامد $A$ باشد، می‌گوییم \emph{$A$ رخ داده است}.
\end{definition}

\begin{example}[label=hs:ex:events]{پیشامدهایی در پرتاب تاس}
در پرتاب یک تاس ($S=\{1,2,3,4,5,6\}$):
\begin{itemize}
  \item «عدد زوج بیاید»: \ $A=\{2,4,6\}$.
  \item «عدد بیشتر از ۴ بیاید»: \ $B=\{5,6\}$.
  \item «عدد ۷ بیاید»: \ $\varnothing$ (مجموعهٔ تهی)؛ این پیشامد \emph{غیرممکن} است.
  \item «عددی بین ۱ تا ۶ بیاید»: \ $S$؛ این پیشامد \emph{حتمی} است.
\end{itemize}
اگر تاس ۶ بیاید، هم $A$ و هم $B$ رخ داده‌اند. اگر ۲ بیاید، $A$ رخ داده ولی $B$ نه.
\end{example}

\subsection{ترکیب پیشامدها: «یا»، «و»، «نه»}

پیشامدها مجموعه‌اند، پس با عمل‌های مجموعه‌ای ترکیبشان می‌کنیم:

\begin{center}
\renewcommand{\arraystretch}{1.4}
\begin{tabular}{>{\semifont}l c l}
\toprule
به زبان عادی & نماد & یعنی \\
\midrule
«$A$ \emph{یا} $B$» (دست‌کم یکی) & $A\cup B$ & اجتماع \\
«$A$ \emph{و} $B$» (هر دو) & $A\cap B$ & اشتراک \\
«$A$ رخ \emph{ندهد}» & $A'$ & متمم (همهٔ برآمدهایی که در $A$ نیستند)\\
\bottomrule
\end{tabular}
\end{center}

\begin{figblock}
\begin{tikzpicture}[scale=0.62]
  \def\univ{(-1.7,-1.35) rectangle (2.9,1.35)}
  \def\sA{(0,0) circle (1)}
  \def\sB{(1.2,0) circle (1)}
  \tikzset{venn/.style={draw=primary, thick}}
  \begin{scope}[xshift=10.8cm]
    \fill[vennfill] \sA; \fill[vennfill] \sB; \draw[venn] \sA \sB; \draw[graytext] \univ;
    \node at (-0.4,0.25) {\lr{$A$}}; \node at (1.6,0.25) {\lr{$B$}};
    \node[font=\small] at (0.6,-1.85) {\lr{$A\cup B$}};
  \end{scope}
  \begin{scope}[xshift=5.4cm]
    \begin{scope}\clip \sA; \fill[vennfill] \sB; \end{scope}
    \draw[venn] \sA \sB; \draw[graytext] \univ;
    \node at (-0.4,0.25) {\lr{$A$}}; \node at (1.6,0.25) {\lr{$B$}};
    \node[font=\small] at (0.6,-1.85) {\lr{$A\cap B$}};
  \end{scope}
  \begin{scope}
    \fill[vennfill] \univ; \fill[white] \sA; \draw[venn] \sA; \draw[graytext] \univ;
    \node at (0,0) {\lr{$A$}};
    \node[font=\small] at (0.6,-1.85) {\lr{$A'$}};
  \end{scope}
\end{tikzpicture}

\figcaption[hs:fig:venn]{نمودار وِن: مستطیل فضای نمونه $S$ است و ناحیهٔ رنگی، پیشامد حاصل.}
\end{figblock}

\begin{definition}[label=hs:def:disjoint]{پیشامدهای ناسازگار}
دو پیشامد $A$ و $B$ را \term{ناسازگار}{mutually exclusive} می‌گوییم اگر نتوانند با هم رخ دهند؛ یعنی $A\cap B=\varnothing$. مثلاً در پرتاب تاس، «زوج آمدن» و «فرد آمدن» ناسازگارند.
\end{definition}

\begin{example}[label=hs:ex:combine]{ترکیب پیشامدها}
در پرتاب تاس، $A=\{2,4,6\}$ (زوج) و $B=\{4,5,6\}$ (بیشتر از ۳). پیشامدهای $A\cup B$، $A\cap B$ و $A'$ را بنویسید و به زبان عادی بیان کنید.
\begin{solution}
\begin{itemize}
  \item $A\cup B=\{2,4,5,6\}$: «زوج یا بیشتر از ۳».
  \item $A\cap B=\{4,6\}$: «زوج و بیشتر از ۳».
  \item $A'=\{1,3,5\}$: «زوج نباشد»، یعنی «فرد».
\end{itemize}
\end{solution}
\end{example}

%% ------------------------------------------------------------
\section{احتمال تجربی: با تکرار آزمایش}
%% ------------------------------------------------------------

یک راه پیدا کردن احتمال این است که آزمایش را \emph{بارها} تکرار کنیم و ببینیم پیشامد مورد نظر چند بار رخ می‌دهد.

\begin{definition}[label=hs:def:empirical]{احتمال تجربی (فراوانی نسبی)}
اگر آزمایشی را $N$ بار تکرار کنیم و پیشامد $A$ در $N_A$ بار رخ دهد، کسر $\dfrac{N_A}{N}$ را \term{فراوانی نسبی}{relative frequency} یا \emph{احتمال تجربی} $A$ می‌نامیم.
\end{definition}

\begin{example}[label=hs:ex:coins-history]{آزمایش‌های واقعی پرتاب سکه}
چند دانشمند در گذشته واقعاً سکه را هزاران بار پرتاب کرده‌اند:
\begin{center}
\renewcommand{\arraystretch}{1.3}
\begin{tabular}{lccc}
\toprule
آزمایشگر & تعداد پرتاب & تعداد شیر & فراوانی نسبی \\
\midrule
بوفون & $4040$ & $2048$ & $0.507$ \\
پیرسون & $24000$ & $12012$ & $0.5005$ \\
\bottomrule
\end{tabular}
\end{center}
هر چه تعداد پرتاب‌ها بیشتر شود، فراوانی نسبی به $\tfrac12$ نزدیک‌تر می‌شود. این واقعیت را \emph{قانون اعداد بزرگ} می‌نامند.
\end{example}

\begin{example}[label=hs:ex:pin]{پونز}
وقتی یک پونز را روی میز می‌اندازیم، یا «نوک به بالا» می‌افتد یا «به پهلو». احتمال «نوک به بالا» چقدر است؟
\begin{solution}
اینجا دو حالت \emph{هم‌شانس نیستند} (شکل پونز متقارن نیست)، پس نمی‌توانیم بگوییم $\tfrac12$. تنها راه، آزمایش است: مثلاً اگر پونز را ۲۰۰ بار بیندازیم و ۱۲۴ بار نوک به بالا بیفتد، احتمال تجربی $\tfrac{124}{200}=0.62$ است. با تکرار بیشتر، تخمین دقیق‌تری به دست می‌آید.
\end{solution}
\end{example}

\begin{warning}
\textbf{خطای قمارباز:} اگر سکهٔ سالمی ۵ بار پشت سر هم شیر بیاید، احتمال اینکه دفعهٔ ششم خط بیاید همچنان $\tfrac12$ است، نه بیشتر! سکه «حافظه» ندارد و نمی‌داند قبلاً چه آمده است. قانون اعداد بزرگ دربارهٔ تعداد \emph{خیلی زیاد} پرتاب است، نه اینکه نتیجه‌ها همدیگر را «جبران» کنند.
\end{warning}

در فصل ۳ راه دوم را می‌بینیم: وقتی برآمدها \emph{هم‌شانس} باشند، احتمال را بدون آزمایش و فقط با \emph{شمردن} حساب می‌کنیم. برای همین، فصل بعد را به شمارش اختصاص داده‌ایم.

%% ------------------------------------------------------------
\section*{خلاصهٔ فصل ۱}
\addcontentsline{toc}{section}{خلاصه و تمرین‌ها}
%% ------------------------------------------------------------

\begin{summary}[title={\leavevmode\boxicon{accent}{$\Sigma$}\ \ خلاصهٔ فصل}]
\begin{itemize}
  \item احتمال عددی بین ۰ (غیرممکن) و ۱ (حتمی) است.
  \item \textbf{فضای نمونه} $S$ مجموعهٔ همهٔ برآمدهای ممکن است و هر زیرمجموعهٔ آن یک \textbf{پیشامد} است.
  \item «یا» $=\cup$، \ «و» $=\cap$، \ «نه» $=$ متمم $A'$. پیشامدهای \textbf{ناسازگار} با هم رخ نمی‌دهند.
  \item \textbf{احتمال تجربی} $=\dfrac{N_A}{N}$؛ با تکرار زیاد به احتمال واقعی نزدیک می‌شود.
  \item سکه و تاس حافظه ندارند (خطای قمارباز).
\end{itemize}
\end{summary}

\subsection*{تمرین‌ها}

\begin{exercise}
فضای نمونهٔ آزمایش‌های زیر را بنویسید و $n(S)$ را بیابید: (الف) پرتاب یک سکه و یک تاس؛ (ب) پرتاب سه سکه؛ (ج) خانواده‌ای با دو فرزند (دختر یا پسر، به ترتیب سن).
\answer{(الف) ۱۲ \quad (ب) ۸ \quad (ج) ۴}
\end{exercise}

\begin{exercise}
در پرتاب تاس، $A$ پیشامد «عدد اول» و $B$ پیشامد «عدد کمتر از ۴» است. $A$، $B$، $A\cup B$، $A\cap B$ و $B'$ را بنویسید.
\answer{$A=\{2,3,5\}$، \ $B=\{1,2,3\}$، \ $A\cup B=\{1,2,3,5\}$، \ $A\cap B=\{2,3\}$، \ $B'=\{4,5,6\}$}
\end{exercise}

\begin{exercise}
کدام زوج‌ها ناسازگارند؟ (الف) «سکه شیر بیاید» و «سکه خط بیاید»؛ (ب) «تاس زوج بیاید» و «تاس ۶ بیاید»؛ (ج) «تاس کمتر از ۳ بیاید» و «تاس بیشتر از ۴ بیاید».
\answer{(الف) و (ج)}
\end{exercise}

\begin{exercise}
یک بسکتبالیست در ۸۰ پرتاب آزاد، ۶۰ بار موفق شده است. احتمال تجربی موفقیت او چقدر است؟ در ۲۰۰ پرتاب بعدی انتظار دارید حدوداً چند بار موفق شود؟
\answer{$0.75$؛ حدود ۱۵۰ بار}
\end{exercise}

\begin{exercise}
علی می‌گوید: «در قرعه‌کشی دیروز عدد ۷ آمد، پس امروز احتمال آمدن ۷ کمتر است.» نظر او درست است؟ توضیح دهید.
\end{exercise}
